\documentclass[12pt,a4paper]{article}
\usepackage[utf8]{inputenc}
\usepackage[T1]{fontenc}
\usepackage[spanish]{babel}
\usepackage{lmodern}
\usepackage{geometry}
\usepackage{enumitem}
\usepackage{float}
\usepackage{pgfplots}
\pgfplotsset{compat=1.18}
\usepackage{hyperref}
\usepackage{amsmath}
\usepackage{graphicx}
\usepackage{xurl}
\usepackage{bm}
\usepackage{booktabs}
\usepackage{multirow}
\usepackage{tikz}
\usetikzlibrary{positioning}
\geometry{left=3cm, right=2.5cm, top=2.5cm, bottom=2.5cm}
\setlength{\parskip}{1em}
\setlength{\parindent}{0em}

\begin{document}

%--------------------------------------------------------------
% PORTADA
%--------------------------------------------------------------
\begin{titlepage}
    \centering
    \vspace*{1cm}
    {\large \textbf{UNIVERSIDAD POPULAR DEL CESAR}}\par
    \vspace{0.3cm}
    {\large FACULTAD DE INGENIERÍA{}S Y TECNOLOGÍAS}\par
    \vspace{3cm}
    {\LARGE \textbf{TALLER PRÁCTICO DE MODELACIÓN:\\[0.4em]
    IDENTIFICACIÓN DE SISTEMAS Y CONDICIONES DE FRONTERA}}\par
    \vspace{3cm}
    {\large \textbf{Autores:}}\par
    \vspace{0.3cm}
    {\large Jorman David Noriega Dávila}\par
    {\large Julio Cesar Ríos Maldonado}\par
    \vspace{3cm}
    {\large \textbf{Programa:} INGENIERÍA DE SISTEMAS}\par
    \vspace{0.3cm}
    {\large \textbf{Asignatura:} MODELOS Y SIMULACIÓN}\par
    \vspace{3cm}
    {\large MARZO DE 2026}\par
    \vspace{0.3cm}
    {\large VALLEDUPAR}\par
\end{titlepage}

\tableofcontents
\newpage

%==============================================================
\section*{CASO DE ESTUDIO 1: TERMODINÁMICA DE UN EXTRUSOR 3D}
\addcontentsline{toc}{section}{Caso de Estudio 1: Termodinámica de un Extrusor 3D}
%==============================================================

\noindent\textbf{Contexto:} Sistema de control de temperatura para el hotend de una impresora FDM.
Filamento de PLA sólido a temperatura ambiente ingresa por la parte superior de un tubo metálico.
Un bloque calentador de aluminio (con cartucho eléctrico) funde el plástico, el cual sale por una
boquilla de latón de 0.4~mm. Un ventilador sopla aire continuamente sobre las aletas superiores
del tubo para evitar fusión prematura.

%--------------------------------------------------------------
\subsection*{Pregunta 1 -- Fronteras y Sistema}
\addcontentsline{toc}{subsection}{Pregunta 1 -- Fronteras y Sistema}
%--------------------------------------------------------------

\noindent El sistema se define como la masa de plástico (filamento PLA) contenida dentro del tubo
metálico, desde el punto de entrada superior hasta la boquilla de salida inferior.
Las fronteras son:

\begin{itemize}
    \item \textbf{Frontera superior (entrada):} Plano transversal donde el filamento PLA sólido a
          temperatura ambiente (${\sim}25$~°C) ingresa al tubo. Es una frontera de tipo
          \textit{inlet} con temperatura conocida.

    \item \textbf{Frontera inferior (salida/boquilla):} Plano donde el plástico fundido sale por la
          boquilla de 0.4~mm. Aquí el plástico abandona el sistema.

    \item \textbf{Frontera lateral:} La interfaz entre el plástico y la pared interior del tubo
          metálico. A través de esta frontera se produce la transferencia de calor (del bloque
          calentador hacia el plástico en la parte baja, y extracción de calor en la parte alta por
          el ventilador).

    \item El metal (tubo, bloque calentador, aletas) y el aire impulsado por el ventilador son
          \textbf{entorno}, no sistema.
\end{itemize}

\noindent A continuación se muestra un diagrama esquemático del hotend con las fronteras
identificadas:

\begin{figure}[H]
\centering
\begin{tikzpicture}[>=latex, thick, font=\small]

    % --- tubo metálico ---
    \fill[gray!25] (-0.6,4) rectangle (0.6,0.5);
    \draw[gray!60, line width=0.8pt] (-0.6,4) rectangle (0.6,0.5);

    % --- bloque calentador ---
    \fill[orange!40] (-0.8,0.5) rectangle (0.8,-0.5);
    \draw[orange!70, line width=0.8pt] (-0.8,0.5) rectangle (0.8,-0.5);
    \node at (0,0) {Bloque calentador};

    % --- boquilla ---
    \fill[yellow!50] (-0.25,-0.5) -- (0.25,-0.5) -- (0.15,-1) -- (-0.15,-1) -- cycle;
    \draw[yellow!80!black, line width=0.8pt]
        (-0.25,-0.5) -- (0.25,-0.5) -- (0.15,-1) -- (-0.15,-1) -- cycle;
    \node[right=0.4cm] at (0.25,-0.75) {Boquilla 0.4~mm};

    % --- filamento entrante ---
    \draw[->, blue!70!black, line width=1.2pt] (0,5.2) -- (0,4.05);
    \node[right=0.2cm] at (0,5.2) {\textbf{Filamento PLA}};
    \node[right=0.2cm] at (0,4.75) {$T_{amb}\approx25$~°C};

    % --- plástico fundido saliente ---
    \draw[->, red!70!black, line width=1.2pt] (0,-1.05) -- (0,-1.7);
    \node[right=0.15cm] at (0,-1.4) {Plástico fundido};

    % --- ventilador / convección lateral ---
    \draw[->, cyan!70!black, line width=1pt] (-1.8,3.5) -- (-0.62,3.5);
    \draw[->, cyan!70!black, line width=1pt] (-1.8,3.0) -- (-0.62,3.0);
    \draw[->, cyan!70!black, line width=1pt] (-1.8,2.5) -- (-0.62,2.5);
    \node[left=0.1cm] at (-1.8,3.0) {Ventilador};

    % --- etiquetas de fronteras ---
    % Frontera superior
    \draw[dashed, red, line width=0.8pt] (-0.8,4) -- (0.8,4);
    \node[right=0.15cm, red] at (0.8,4) {\textbf{Frontera superior}};
    \node[right=0.15cm, red] at (0.8,3.65) {(Dirichlet: $T=T_{amb}$)};

    % Frontera lateral
    \draw[dashed, blue!70, line width=0.8pt] (0.62,3.0) -- (1.8,3.0);
    \node[right=0.1cm, blue!70] at (1.8,3.1) {\textbf{Frontera lateral}};
    \node[right=0.1cm, blue!70] at (1.8,2.75) {(Robin: convección)};

    % Frontera inferior
    \draw[dashed, green!60!black, line width=0.8pt] (-0.85,-0.5) -- (0.85,-0.5);
    \node[right=0.15cm, green!60!black] at (0.85,-0.5) {\textbf{Frontera inferior}};
    \node[right=0.15cm, green!60!black] at (0.85,-0.8) {(Salida/boquilla)};

    % Eje axial
    \draw[dotted, gray] (0,4) -- (0,0.5);
    \node[left=0.05cm, gray] at (0,2.25) {$x$};
    \draw[->, gray] (0,2.5) -- (0,1.8);

\end{tikzpicture}
\caption{Diagrama esquemático del hotend FDM con fronteras del sistema identificadas.}
\label{fig:hotend}
\end{figure}

%--------------------------------------------------------------
\subsection*{Pregunta 2 -- Parámetros vs.\ Variables}
\addcontentsline{toc}{subsection}{Pregunta 2 -- Parámetros vs. Variables}
%--------------------------------------------------------------

\noindent\textbf{Parámetros físicos:}
\begin{enumerate}
    \item \textbf{Conductividad térmica del PLA:} $k_{PLA}$ [W/(m·K)] --- constante que
          caracteriza cuánto calor conduce el material.
    \item \textbf{Calor específico del PLA:} $c_p$ [J/(kg·K)] --- determina cuánta energía se
          necesita para elevar la temperatura del material.
    \item \textbf{Densidad del PLA:} $\rho$ [kg/m³].
    \item \textbf{Potencia del cartucho calentador:} $\dot{Q}_{calentador}$ [W] --- parámetro de
          entrada al sistema.
\end{enumerate}

\noindent\textbf{Variable de estado principal:} La temperatura del plástico $T(x, t)$ [°C o K],
que varía a lo largo del eje axial del tubo ($x$) y en el tiempo ($t$).

%--------------------------------------------------------------
\subsection*{Pregunta 3 -- Condiciones de Frontera}
\addcontentsline{toc}{subsection}{Pregunta 3 -- Condiciones de Frontera}
%--------------------------------------------------------------

\noindent\textbf{Frontera de entrada (plástico frío):}

En la entrada superior, la temperatura del plástico es conocida e igual a la temperatura ambiente.
Esta es una condición de \textbf{Dirichlet} (condición de primer tipo), porque se prescribe
directamente el valor de la variable dependiente (temperatura).

\begin{equation}
    T\big|_{x=0} = T_{amb} = 25\,\text{°C} \qquad \text{(Condición de Dirichlet)}
\end{equation}

\noindent\textbf{Frontera lateral (metal enfriado por ventilador):}

En la frontera lateral del metal donde actúa el ventilador, el flujo de calor depende de la
diferencia entre la temperatura de la superficie y la del aire. Esto corresponde a una condición
de \textbf{Robin} (condición de tercer tipo o convectiva):

\begin{equation}
    -k \frac{\partial T}{\partial n}\bigg|_{\text{lateral}} = h\,(T_{sup} - T_{\infty})
\end{equation}

\noindent donde $h$ es el coeficiente de convección forzada del ventilador, $T_{sup}$ es la
temperatura de la superficie del metal y $T_{\infty}$ es la temperatura del aire. No se prescribe
el valor de $T$ ni el flujo por separado, sino una combinación lineal de ambos.

%--------------------------------------------------------------
\subsection*{Pregunta 4 -- Decisión de Modelado (¿Parámetros concentrados?)}
\addcontentsline{toc}{subsection}{Pregunta 4 -- Decisión de Modelado}
%--------------------------------------------------------------

\noindent\textbf{No se puede modelar el plástico dentro del tubo como sistema de parámetros
concentrados.} Justificación:

\begin{itemize}
    \item \textbf{Criterio físico:} El tubo tiene una longitud considerable (${\sim}30$--$40$~mm en
          hotends típicos) y el plástico experimenta un gradiente de temperatura fuerte y no uniforme
          a lo largo del eje axial: a la entrada está a ${\sim}25$~°C (sólido) y en la punta
          (boquilla) debe estar a ${\sim}210$~°C (fundido). La temperatura \textbf{NO} es homogénea
          en el espacio.

    \item \textbf{Criterio matemático (Número de Biot):} Para que un sistema de parámetros
          concentrados sea válido, se requiere:
          \begin{equation}
              Bi = \frac{h\,L_c}{k} \ll 0.1
          \end{equation}
          El PLA tiene baja conductividad térmica ($k \approx 0.13$~W/(m·K)), lo que implica un
          $Bi$ elevado, indicando que los gradientes espaciales internos son significativos y
          \textbf{NO} pueden despreciarse.

    \item \textbf{Conclusión:} Se requiere un modelo de \textbf{parámetros distribuidos}, descrito
          por una Ecuación Diferencial Parcial (EDP) de calor unidimensional:
          \begin{equation}
              \rho\,c_p\,\frac{\partial T}{\partial t} = k\,\frac{\partial^2 T}{\partial x^2}
              + \dot{q}(x,t)
          \end{equation}
\end{itemize}

\newpage

%==============================================================
\section*{CASO DE ESTUDIO 2: INSTALACIÓN DE UN ARREGLO DE PANELES SOLARES}
\addcontentsline{toc}{section}{Caso de Estudio 2: Instalación de un Arreglo de Paneles Solares}
%==============================================================

\noindent\textbf{Contexto:} Sistema fotovoltaico sobre techo en zona de alta temperatura y
radiación. El panel convierte parte de la luz en electricidad y otra parte en calor residual.
Transfiere calor al aire de la calle (cara superior, expuesta al viento) y al aire estancado entre
panel y tejas (cara inferior). La energía eléctrica generada fluye hacia un banco de baterías de
litio dentro de la casa.

%--------------------------------------------------------------
\subsection*{Pregunta 1 -- Sistemas Múltiples}
\addcontentsline{toc}{subsection}{Pregunta 1 -- Sistemas Múltiples}
%--------------------------------------------------------------

\noindent\textbf{Sistema 1 -- Dominio Térmico (el panel solar):}

\begin{itemize}
    \item \textbf{Sistema:} El panel fotovoltaico (celda + encapsulante + vidrio).
    \item \textbf{Entradas:} Irradiancia solar $G$ [W/m²], temperatura del aire exterior $T_{amb}$,
          velocidad del viento.
    \item \textbf{Salidas:} Temperatura del panel $T_{panel}$, potencia eléctrica generada
          $P_{elec}$ (que sale del dominio térmico hacia el dominio eléctrico).
    \item \textbf{Ecuación gobernante} (balance de energía térmica):
          \begin{equation}
              G \cdot \alpha = P_{elec} + Q_{conv\_sup} + Q_{conv\_inf}
          \end{equation}
\end{itemize}

\noindent\textbf{Sistema 2 -- Dominio Eléctrico (banco de baterías):}

\begin{itemize}
    \item \textbf{Sistema:} El banco de baterías de litio y el circuito de carga.
    \item \textbf{Entradas:} Potencia eléctrica generada por el panel $P_{elec}(t)$ [W].
    \item \textbf{Salidas:} Estado de Carga (SOC) de las baterías, voltaje y corriente del banco.
    \item \textbf{Ecuación gobernante:}
          \begin{equation}
              \frac{d(\text{SOC})}{dt} = \frac{I(t)}{C_{\text{batería}}}
          \end{equation}
\end{itemize}

\noindent\textbf{Interconexión:} La salida eléctrica del Sistema~1 ($P_{elec}$) es la entrada del
Sistema~2. La temperatura del panel (Sistema~1) afecta la eficiencia y por tanto $P_{elec}$,
creando un acoplamiento bidireccional.

%--------------------------------------------------------------
\subsection*{Pregunta 2 -- Condiciones de Frontera Térmicas (Radiación Solar)}
\addcontentsline{toc}{subsection}{Pregunta 2 -- Condiciones de Frontera Térmicas}
%--------------------------------------------------------------

\noindent La radiación solar de 1000~W/m² representa una condición de \textbf{Neumann}
(condición de segundo tipo).

\noindent\textbf{Justificación:} Una condición de Neumann prescribe el \textbf{flujo} (derivada
normal) de la variable dependiente en la frontera, no su valor. En este caso, se conoce el flujo
de calor entrante por unidad de área:

\begin{equation}
    -k\frac{\partial T}{\partial n}\bigg|_{\text{cara superior}} = q''_{solar}
    = 1000\,\text{W/m}^2
\end{equation}

\noindent No se prescribe la temperatura (eso sería Dirichlet), sino la cantidad de energía que
atraviesa la frontera por unidad de tiempo y área. Por eso es condición de \textbf{Neumann}.

%--------------------------------------------------------------
\subsection*{Pregunta 3 -- La Condición de Robin (cara superior con viento)}
\addcontentsline{toc}{subsection}{Pregunta 3 -- La Condición de Robin}
%--------------------------------------------------------------

\noindent La cara superior del panel, barrida por ráfagas de viento, exige una condición de
\textbf{Robin} porque simultáneamente ocurren dos fenómenos en esa misma frontera:

\begin{enumerate}
    \item El panel recibe flujo de calor por radiación (Neumann parcial).
    \item El panel intercambia calor con el aire por \textbf{convección forzada} (dependiente de
          la temperatura).
\end{enumerate}

\noindent La condición de Robin combina la temperatura de la superficie y el gradiente en la
frontera:

\begin{equation}
    -k\frac{\partial T}{\partial n}\bigg|_{sup} = h_{conv}\,(T_{sup} - T_{\infty})
\end{equation}

\noindent\textbf{Variables que interactúan:}
\begin{itemize}
    \item $T_{sup}$: temperatura de la superficie del panel [°C o K].
    \item $T_{\infty}$: temperatura del aire ambiente [°C o K].
    \item $h_{conv}$: coeficiente de convección forzada por el viento [W/(m²·K)] --- variable
          porque depende de la velocidad del viento.
    \item $k$: conductividad térmica del panel [W/(m·K)].
    \item $\dfrac{\partial T}{\partial n}$: gradiente de temperatura en dirección normal a la
          superficie.
\end{itemize}

\noindent La condición no puede ser solo Dirichlet (no se conoce $T_{sup}$) ni solo Neumann (el
flujo depende de $T_{sup}$), sino la combinación \textbf{Robin}.

%--------------------------------------------------------------
\subsection*{Pregunta 4 -- Distribuidos vs.\ Concentrados}
\addcontentsline{toc}{subsection}{Pregunta 4 -- Distribuidos vs. Concentrados}
%--------------------------------------------------------------

\noindent\textbf{Temperatura de una celda solar a lo largo del día $\rightarrow$ Modelo
CONCENTRADO:}

\begin{itemize}
    \item A lo largo del día, interesa cómo evoluciona la temperatura de la celda \textbf{en el
          tiempo} ($T_{celda}(t)$).
    \item Una celda solar es pequeña (del orden de cm), tiene alta conductividad y el gradiente
          espacial interno es pequeño en comparación con las variaciones temporales.
    \item Se verifica $Bi \ll 0.1$, por lo que la temperatura puede considerarse uniforme en todo
          el volumen de la celda en cada instante.
    \item Modelo resultante (EDO):
          \begin{equation}
              m\,c_p\,\frac{dT_{celda}}{dt} = \dot{Q}_{rad} - \dot{Q}_{conv} - P_{elec}
          \end{equation}
\end{itemize}

\noindent\textbf{SOC del banco de baterías $\rightarrow$ Modelo CONCENTRADO:}

\begin{itemize}
    \item El SOC representa la carga total almacenada en todo el banco como una sola cantidad
          escalar.
    \item No interesa la distribución espacial de la carga dentro de cada batería, sino el nivel
          global de energía.
    \item El banco se modela como un único ``depósito'' de energía:
          \begin{equation}
              \frac{d(\text{SOC})}{dt} = \frac{\eta \cdot I(t)}{C_{nominal}}
          \end{equation}
    \item Modelo resultante: EDO simple $\rightarrow$ parámetros concentrados.
\end{itemize}

\noindent\textbf{Contraste:} Ambos casos usan modelos concentrados, pero por razones distintas:
la celda porque es físicamente pequeña (gradientes espaciales despreciables), y el banco de
baterías porque el SOC es por definición una variable integral/global que no tiene distribución
espacial relevante para el control.

\newpage

%==============================================================
\section*{CASO DE ESTUDIO 3: VIBRACIÓN EN EL ESLABÓN DE UN BRAZO ROBÓTICO}
\addcontentsline{toc}{section}{Caso de Estudio 3: Vibración en el Eslabón de un Brazo Robótico}
%==============================================================

\noindent\textbf{Contexto:} Brazo robótico industrial con eslabón de aluminio de 1.5~m de
longitud. Al frenar bruscamente, la inercia de la pieza pesada en la punta causa vibración y
flexión del brazo, afectando la precisión del ensamblaje.

%--------------------------------------------------------------
\subsection*{Pregunta 1 -- El Sistema y la Frontera}
\addcontentsline{toc}{subsection}{Pregunta 1 -- El Sistema y la Frontera}
%--------------------------------------------------------------

\begin{itemize}
    \item \textbf{Sistema:} El eslabón de aluminio como cuerpo continuo elástico deformable, de
          longitud $L = 1.5$~m.
    \item \textbf{Fronteras físicas:}
          \begin{itemize}
              \item \textit{Extremo base ($x = 0$):} La unión rígida con el motor/actuador. Es
                    una frontera empotrada.
              \item \textit{Extremo punta ($x = L$):} Donde se acopla la pieza o herramienta,
                    que aplica fuerzas y momentos.
              \item \textit{Frontera lateral (superficial):} La superficie exterior del eslabón,
                    en contacto con el ambiente (puede considerarse libre de esfuerzos normales).
          \end{itemize}
    \item \textbf{Entrada de energía mecánica:} El \textbf{torque de frenado} aplicado por el
          motor en la base ($x = 0$) es la entrada de energía mecánica. Cuando el motor frena
          bruscamente, impone una deceleración angular que genera una fuerza de inercia distribuida
          a lo largo del eslabón y una fuerza puntual de la masa en la punta, excitando los modos
          de vibración.
\end{itemize}

%--------------------------------------------------------------
\subsection*{Pregunta 2 -- Parámetros del Eslabón}
\addcontentsline{toc}{subsection}{Pregunta 2 -- Parámetros del Eslabón}
%--------------------------------------------------------------

\noindent Tres parámetros inherentes al material o geometría que determinan la flexión:

\begin{enumerate}
    \item \textbf{Módulo de Elasticidad (Módulo de Young):} $E_{Al} \approx 69$~GPa --- resistencia
          del aluminio a la deformación elástica; mayor $E$ implica menor deflexión.
    \item \textbf{Momento de Inercia de Área de la sección transversal:} $I$~[m\textsuperscript{4}]
          --- depende de la forma geométrica de la sección (circular hueco, rectangular, etc.);
          determina la rigidez a la flexión $EI$.
    \item \textbf{Densidad lineal de masa:} $\rho A$~[kg/m] --- masa por unidad de longitud;
          determina la inercia distribuida del eslabón y las frecuencias naturales de vibración.
\end{enumerate}

%--------------------------------------------------------------
\subsection*{Pregunta 3 -- Condiciones de Frontera Espaciales}
\addcontentsline{toc}{subsection}{Pregunta 3 -- Condiciones de Frontera Espaciales}
%--------------------------------------------------------------

\noindent\textbf{Extremo base (empotrado, $x = 0$):}

El desplazamiento transversal es cero y también la pendiente (rotación) es cero. Esta es una
condición de \textbf{Dirichlet} (primer tipo) --- se prescribe directamente el valor de la
variable dependiente:

\begin{equation}
    v(0,t) = 0 \qquad \text{y} \qquad
    \frac{\partial v}{\partial x}\bigg|_{x=0} = 0 \qquad \text{(Dirichlet)}
\end{equation}

\noindent\textbf{Extremo punta ($x = L$):}

En la punta hay una masa que ejerce una fuerza cortante específica y un momento flector conocido.
Matemáticamente:

\begin{equation}
    EI\frac{\partial^3 v}{\partial x^3}\bigg|_{x=L} = m_{punta}\frac{\partial^2 v}{\partial t^2}
    \bigg|_{x=L}
\end{equation}

Esta es una condición de \textbf{Neumann} (segundo tipo) --- se prescribe la derivada de la
variable (relacionada con el cortante/momento, que son derivadas espaciales del desplazamiento).
No se conoce el desplazamiento en la punta, sino la fuerza/momento que actúa en ella.

%--------------------------------------------------------------
\subsection*{Pregunta 4 -- El Reto Matemático (¿Parámetros concentrados o EDP?)}
\addcontentsline{toc}{subsection}{Pregunta 4 -- El Reto Matemático}
%--------------------------------------------------------------

\noindent\textbf{Rechazamos la suposición del modelo masa-resorte-amortiguador (parámetros
concentrados).}

\noindent\textbf{Justificación física y matemática:}

El eslabón de aluminio de 1.5~m es un \textbf{sólido continuo y distribuido}: en cada punto $x$
a lo largo de su longitud existe un desplazamiento transversal $v(x,t)$ diferente. Cuando el
brazo vibra, no se mueve como un cuerpo rígido; se \textbf{flexiona}, lo que significa que la
deformación varía continuamente en el espacio.

Un modelo masa-resorte-amortiguador (EDO) supone que toda la masa se concentra en un punto y
que la deformación es uniforme --- esto equivale a asumir que el brazo es infinitamente rígido
excepto en un punto, lo cual es físicamente incorrecto para un eslabón largo y esbelto.

\noindent\textbf{¿Por qué la flexión exige una EDP?}

La deflexión $v(x,t)$ depende tanto de la posición $x$ a lo largo del eslabón como del tiempo
$t$. La ecuación que gobierna la vibración transversal de una viga de Euler-Bernoulli es:

\begin{equation}
    \rho A\,\frac{\partial^2 v}{\partial t^2} + EI\,\frac{\partial^4 v}{\partial x^4} = f(x,t)
\end{equation}

\noindent Esta es una \textbf{Ecuación Diferencial Parcial} (EDP) de cuarto orden en el espacio
y segundo orden en el tiempo, con \textbf{dos variables independientes} ($x$ y $t$). Ninguna EDO
puede capturar esta dependencia simultánea.

\noindent\textbf{Consecuencias prácticas:} Un modelo concentrado predice una única frecuencia
natural, mientras que el eslabón real tiene infinitos modos de vibración. Para garantizar
precisión milimétrica en el ensamblaje, es indispensable conocer al menos los primeros modos
flexurales, lo cual solo es posible con la EDP completa o su aproximación por elementos finitos
(FEM).

\end{document}
